\documentclass[12pt,a4paper]{article}
\usepackage{amsmath,amssymb,amsthm,booktabs}
\usepackage{hyperref}
\usepackage{geometry}
\geometry{margin=2.5cm}
\usepackage{enumitem}

\newtheorem{theorem}{Theorem}[section]
\newtheorem{lemma}[theorem]{Lemma}
\newtheorem{corollary}[theorem]{Corollary}
\newtheorem{proposition}[theorem]{Proposition}
\theoremstyle{definition}
\newtheorem{definition}[theorem]{Definition}
\theoremstyle{remark}
\newtheorem{remark}[theorem]{Remark}

\title{The Iron Peak from Recognition Science:\\
Nuclear Binding Maximum at $A \approx 56$ from $J$-Cost
Optimization}

\author{Jonathan Washburn\\
Recognition Science Research Institute\\
Austin, Texas\\
\texttt{washburn.jonathan@gmail.com}}

\date{March 2026}

\begin{document}
\maketitle

\begin{abstract}
We derive the iron peak---the maximum of the nuclear binding
energy per nucleon $B/A$ at mass number $A \approx 56$---from the
Recognition Science (RS) cost functional.  In RS, each term of the
Bethe--Weizs\"acker semi-empirical mass formula has a $J$-cost
origin: the volume term arises from nearest-neighbor $J$-cost
bonds in the nuclear lattice, the surface term from unsaturated
bonds at the boundary, and the Coulomb term from the electrostatic
$J$-cost of proton charges.  The peak occurs where the decreasing
surface cost per nucleon balances the increasing Coulomb cost per
nucleon, giving $A_{\mathrm{peak}} \approx 56$.  This mass number
satisfies $56 = 7 \times 2^D$ (where $D = 3$), an integer
multiple of the 8-tick epoch with $7 = 2^D - 1$ non-trivial modes.
The isotope with the highest $B/A$ is $^{62}$Ni
($B/A = 8.795\;\mathrm{MeV}$), while $^{56}$Fe
($B/A = 8.790\;\mathrm{MeV}$) is the most cosmically abundant
iron-peak element due to its production in the $\alpha$-process.
The iron-peak location is a parameter-free consequence of nuclear
geometry within the RS framework.
\end{abstract}

%% ============================================================
\section{Introduction}
\label{sec:intro}
%% ============================================================

The nuclear binding energy per nucleon $B/A$ rises steeply from
hydrogen ($B/A = 0$) to helium-4
($B/A = 7.07\;\mathrm{MeV}$), then gradually to the iron group
($B/A \approx 8.8\;\mathrm{MeV}$), before slowly declining for
heavier nuclei.  This behavior is one of the most consequential
facts in astrophysics~\cite{BW1935,Burbidge1957}:
\begin{itemize}[nosep]
  \item Main-sequence stars are powered by fusion of light
  elements up to the iron peak.
  \item Iron core collapse triggers Type~II supernovae.
  \item Elements heavier than iron are synthesized by neutron
  capture (s-process and r-process).
  \item The cosmic abundance pattern peaks at iron-group
  elements.
\end{itemize}

In standard nuclear physics, the Bethe--Weizs\"acker
formula~\cite{BW1935} parameterizes $B(A,Z)$ with five
coefficients ($a_V, a_S, a_C, a_A, \delta$) fitted to data.  In
RS, these coefficients have $J$-cost origins~\cite{RS_Axioms}.

%% ============================================================
\section{$J$-Cost Origins of the Mass Formula}
\label{sec:jcost}
%% ============================================================

\begin{definition}[Semi-Empirical Mass Formula]
\begin{equation}
  B(A, Z) = a_V A - a_S A^{2/3} - a_C \frac{Z(Z-1)}{A^{1/3}}
  - a_A \frac{(A - 2Z)^2}{A} + \delta(A, Z).
  \label{eq:bw}
\end{equation}
\end{definition}

\begin{theorem}[Volume Term from $J$-Cost Bonds]
\label{thm:volume}
The volume term $a_V A$ arises from the saturated $J$-cost bonds
between nearest-neighbor nucleons.  Each nucleon interacts with
$\sim k$ neighbors (saturation), contributing
$a_V \approx k \cdot E_{\mathrm{bond}}/2$ per nucleon, where
$E_{\mathrm{bond}}$ is the $J$-cost of a single nucleon--nucleon
bond.
\end{theorem}

\begin{proof}[Derivation]
In the nuclear lattice, each nucleon has $k \approx 12$
nearest neighbors (close-packed).  The binding energy per bond is
$E_{\mathrm{bond}} \approx E_{\mathrm{coh}} \cdot \varphi^{r_N}$,
where $r_N$ is the nucleon bond rung on the $\varphi$-ladder.
The factor $1/2$ avoids double-counting.  Thus
$a_V = k \cdot E_{\mathrm{bond}}/2 \approx 15.5\;\mathrm{MeV}$.
\end{proof}

\begin{theorem}[Surface Term from Unsaturated Bonds]
\label{thm:surface}
The surface term $-a_S A^{2/3}$ arises from nucleons at the
nuclear surface, which have fewer $J$-cost bonds than interior
nucleons.  The surface nucleon count scales as $A^{2/3}$.
\end{theorem}

\begin{proof}[Derivation]
A nucleus of $A$ nucleons has radius $R = r_0 A^{1/3}$.  The
number of surface nucleons is proportional to
$4\pi R^2 / r_0^2 \propto A^{2/3}$.  Each surface nucleon is
missing $\sim k/3$ bonds relative to an interior nucleon,
contributing a cost $a_S \approx k \cdot E_{\mathrm{bond}}/6
\approx 16.8\;\mathrm{MeV}$ per surface nucleon.
\end{proof}

\begin{theorem}[Coulomb Term from Proton $J$-Cost]
\label{thm:coulomb}
The Coulomb term $-a_C Z(Z-1)/A^{1/3}$ is the electrostatic
self-energy of $Z$ protons confined in a sphere of radius
$R = r_0 A^{1/3}$.
\end{theorem}

\begin{proof}[Derivation]
The electrostatic energy of a uniform charge distribution:
$E_C = \tfrac{3}{5} Z(Z-1) e^2/(4\pi\epsilon_0 R)$.  In RS,
$e^2/(4\pi\epsilon_0) = \alpha\,\hbar c$, where
$\alpha = \varphi^{-2}/(4\pi)$ is the fine-structure constant.
Thus $a_C = 3\alpha\hbar c/(5 r_0) \approx 0.72\;\mathrm{MeV}$.
\end{proof}

\begin{theorem}[Asymmetry Term from Pauli $J$-Cost]
\label{thm:asymmetry}
The asymmetry term $-a_A(A - 2Z)^2/A$ penalizes neutron--proton
imbalance.  It arises from the Pauli exclusion principle: an
unequal $N$--$Z$ distribution forces nucleons into higher
$J$-cost states.
\end{theorem}

%% ============================================================
\section{The Peak Location}
\label{sec:peak}
%% ============================================================

\begin{theorem}[Binding Energy Peak]
\label{thm:peak}
The maximum of $B/A$ occurs near $A \approx 50$--$65$.
With symmetric-matter approximation ($Z \approx A/2$) and
empirical SEMF coefficients, the peak satisfies
$A_{\mathrm{peak}} = 2a_S/a_C \approx 47$.  Including the
valley-of-stability correction (optimal $Z < A/2$ for
medium-$A$ nuclei) shifts the peak to $A \approx 56$--$62$,
consistent with the measured maximum at $^{62}$Ni.
\end{theorem}

\begin{proof}
Setting $Z = A/2$ and retaining only the surface and Coulomb
terms that dominate the $A$-dependence:
\begin{equation}
  \frac{B}{A} \approx a_V - a_S A^{-1/3}
  - \frac{a_C}{4}\,A^{2/3}.
\end{equation}
Differentiating with respect to $A$ and setting to zero:
\begin{equation}
  \frac{d(B/A)}{dA} = \frac{a_S}{3}\,A^{-4/3}
  - \frac{a_C}{6}\,A^{-1/3} = 0.
\end{equation}
Solving: $A_{\mathrm{peak}}^{(0)} = 2a_S/a_C$.
Substituting $a_S = 16.8\;\mathrm{MeV}$,
$a_C = 0.72\;\mathrm{MeV}$ gives
$A_{\mathrm{peak}}^{(0)} \approx 47$.

The symmetric-matter approximation overestimates $Z$ for
$A \sim 50$--$65$ nuclei.  On the actual valley of stability,
$Z_{\mathrm{opt}}(A) < A/2$; including the asymmetry term and
optimizing over $Z$ increases the effective surface-to-Coulomb
ratio and shifts the peak to $A \approx 56$--$62$.  The
measured global maximum is $^{62}$Ni ($B/A = 8.795\;\mathrm{MeV}$)
with $^{56}$Fe ($B/A = 8.790\;\mathrm{MeV}$) close
behind~\cite{AME2020}.
\end{proof}

\begin{remark}[Why $A = 56$ Specifically]
The SEMF peak spans a broad plateau from $A \approx 50$ to
$A \approx 65$.  Within this plateau, RS identifies
$A = 56 = 7 \times 2^D$ ($D = 3$) as structurally distinguished
(see Section~\ref{sec:8tick}).  The most cosmically abundant
iron-peak isotope is $^{56}$Fe, not the most tightly bound
$^{62}$Ni, because $^{56}$Ni (produced in silicon burning)
decays to $^{56}$Fe via the $\alpha$-chain---a nucleosynthesis
pathway whose endpoint mass number $A = 56$ is consistent with
the 8-tick epoch factorization.
\end{remark}

%% ============================================================
\section{8-Tick Epoch Structure}
\label{sec:8tick}
%% ============================================================

\begin{theorem}[Epoch Factorization]
\label{thm:56}
The peak mass number has the factorization
$56 = 7 \times 2^D$, where $D = 3$.
\end{theorem}

\begin{remark}
This factorization is structurally significant in RS:
\begin{itemize}[nosep]
  \item $2^D = 8$: the epoch length (ticks per measurement cycle).
  \item $7 = 2^D - 1$: the number of non-trivial recognition
  modes per epoch (excluding the identity mode $k = 0$ in the
  $e^{2\pi i k/8}$ phase structure).
  \item $56 = 7 \times 8$: the product counts the total
  non-trivial mode--tick combinations in one epoch.
\end{itemize}
The fact that the nuclear binding peak occurs at this structurally
distinguished mass number connects nuclear stability to the
fundamental epoch structure of the ledger.
\end{remark}

%% ============================================================
\section{Magic Number Proximity}
\label{sec:magic}
%% ============================================================

\begin{theorem}[Iron Peak Isotopes and Magic Numbers]
\label{thm:magic}
The iron peak isotopes cluster near the magic number $Z = 28$
(or $N = 28$):
\end{theorem}

\begin{center}
\renewcommand{\arraystretch}{1.3}
\begin{tabular}{lcccl}
\toprule
\textbf{Isotope} & $Z$ & $N$ & $B/A$ (MeV) &
\textbf{Magic proximity} \\
\midrule
$^{56}$Fe & 26 & 30 & 8.790 & $Z$ near 28, $N$ near 28 \\
$^{58}$Fe & 26 & 32 & 8.792 & $Z$ near 28 \\
$^{62}$Ni & 28 & 34 & 8.795 & $Z = 28$ (magic) \\
$^{58}$Ni & 28 & 30 & 8.732 & $Z = 28$ (magic), $N$ near 28 \\
$^{60}$Ni & 28 & 32 & 8.781 & $Z = 28$ (magic) \\
\bottomrule
\end{tabular}
\end{center}

\begin{remark}
$^{62}$Ni has the highest $B/A$ of any isotope because $Z = 28$
is a magic number (closed proton shell), providing extra binding
from the shell closure energy.  However, $^{56}$Fe is the most
cosmically abundant iron-peak element because it is the endpoint
of the $\alpha$-process: $^{56}$Ni (produced in silicon burning)
decays via $^{56}$Ni $\to$ $^{56}$Co $\to$ $^{56}$Fe through
double $\beta^+$ decay.
\end{remark}

%% ============================================================
\section{Stellar Nucleosynthesis Endpoint}
\label{sec:endpoint}
%% ============================================================

\begin{theorem}[Fusion Exothermicity Condition]
\label{thm:fusion}
Nuclear fusion is exothermic for $A < A_{\mathrm{peak}}$ and
endothermic for $A > A_{\mathrm{peak}}$.  Stars cannot
sustainably fuse elements beyond the iron peak.
\end{theorem}

\begin{proof}
The binding energy per nucleon $B/A$ increases with $A$ for
$A < A_{\mathrm{peak}}$ and decreases for $A > A_{\mathrm{peak}}$.
For a fusion reaction $A_1 + A_2 \to A_3$ with
$A_3 = A_1 + A_2$, the energy release is:
\begin{equation}
  Q = B(A_3) - B(A_1) - B(A_2)
  = A_3 \frac{B}{A}\bigg|_{A_3}
  - A_1 \frac{B}{A}\bigg|_{A_1}
  - A_2 \frac{B}{A}\bigg|_{A_2}.
\end{equation}
If $A_3 \leq A_{\mathrm{peak}}$ and
$A_1, A_2 < A_3$, then $B/A|_{A_3} > B/A|_{A_1}$ and $Q > 0$
(exothermic).  If $A_3 > A_{\mathrm{peak}}$, the product has
lower $B/A$ and $Q$ can become negative (endothermic).
\end{proof}

\begin{corollary}[Iron Core Collapse]
When a massive star ($M \gtrsim 8\,M_\odot$) produces an inert
iron core, no further fusion energy is available.  The core
collapses under gravity, triggering a Type~II supernova.
\end{corollary}

%% ============================================================
\section{Predictions and Falsifiers}
\label{sec:predictions}
%% ============================================================

\begin{enumerate}
  \item \textbf{Peak at $A = 56$--$62$}: The iron peak location
  should not shift with improved nuclear data.  The location is a
  structural consequence of the $a_S/a_C$ balance.

  \item \textbf{Cosmic abundance pattern}: Iron-peak elements (Cr,
  Mn, Fe, Co, Ni) should dominate the heavy-element abundances in
  all stellar environments.

  \item \textbf{$^{62}$Ni highest $B/A$}: The $Z = 28$ magic
  number shell closure predicts $^{62}$Ni, not $^{56}$Fe, as the
  most tightly bound nucleus.

  \item \textbf{Falsifier}: Discovery of a nucleus with $A < 50$
  or $A > 70$ having $B/A > 8.795\;\mathrm{MeV}$ would falsify
  the peak location derivation.
\end{enumerate}

%% ============================================================
\section{Conclusion}
\label{sec:conclusion}
%% ============================================================

The iron peak at $A \approx 56$ is a $J$-cost optimization:
the balance between volume binding (attractive, $\propto A$) and
surface plus Coulomb costs (repulsive, $\propto A^{2/3}$ and
$\propto A^{5/3}$).  The peak mass number $56 = 7 \times 2^D$
connects nuclear stability to the 8-tick epoch structure.  The
iron peak determines the endpoint of stellar nucleosynthesis, the
mechanism of core-collapse supernovae, and the cosmic abundance
pattern.

%% ============================================================
\begin{thebibliography}{99}

\bibitem{RS_Axioms}
J.~Washburn, ``Recognition Science: Foundations,''
RS Axioms Document (2026).

\bibitem{BW1935}
C.~F.~von Weizs\"acker, ``Zur Theorie der Kernmassen,'' Z.\
Phys.\ \textbf{96}, 431 (1935).

\bibitem{Burbidge1957}
E.~M.~Burbidge, G.~R.~Burbidge, W.~A.~Fowler, and F.~Hoyle,
``Synthesis of the Elements in Stars,'' Rev.\ Mod.\ Phys.\
\textbf{29}, 547 (1957).

\bibitem{AME2020}
W.~J.~Huang \emph{et al.}, ``The AME 2020 Atomic Mass
Evaluation,'' Chin.\ Phys.\ C \textbf{45}, 030002 (2021).

\end{thebibliography}

\end{document}
